% Partie : sets-functions-relations
% Chapitre : sets
% Section : unions-and-intersections

\documentclass[../../../include/open-logic-section]{subfiles}

\begin{document}

\olfileid[fr]{sfr}{set}{uni}
\olsection{Réunions et intersections}

\begin{explain}
Dans la \olref[sfr][set][bas]{sec}, nous avons introduit les définitions
d'ensembles en compréhension, c'est-à-dire de la forme
$\Setabs{x}{\phi(x)}$. On fait ici appel à une propriété~$\phi$,
qui peut mentionner des ensembles déjà définis. Ainsi, si $A$ et~$B$
sont des ensembles, $\Setabs{x}{x \in A \lor x \in B}$ est formé
de tous les objets qui sont des !!{element}s de $A$ ou de~$B$ :
il rassemble les !!{element}s de $A$ et de~$B$.
La \olref{fig:union} représente cette situation : la zone colorée
correspond aux !!{element}s des deux ensembles $A$ et~$B$ réunis.

\begin{figure}
  \olasset{assets/diagrams/union.tikz}
  \caption{La réunion $A \cup B$ de deux ensembles rassemble
    les !!{element}s de $A$ et ceux de~$B$.}
  \ollabel{fig:union}
\end{figure}

Cette opération qui consiste à réunir des ensembles est très
utile et fréquente ; donnons-lui un nom et un symbole.
\end{explain}

\begin{defn}[Réunion]
La \emph{réunion} de deux ensembles $A$ et $B$, notée $A \cup B$,
est l'ensemble de tous les objets qui sont des !!{element}s
de $A$, de $B$, ou des deux.
\[
A \cup B = \Setabs{x}{x \in A \lor x \in B}
\]
\end{defn}

\begin{ex}
Puisque la répétition des !!{element}s n'importe pas, la réunion de
deux ensembles ayant !!a{element} en commun ne contient cet
!!{element} qu'une seule fois. Par exemple,
$\{ a, b, c\} \cup \{ a, 0, 1\} = \{a, b, c, 0, 1\}$.

La réunion d'un ensemble et de l'un de ses sous-ensembles est
simplement le plus grand des deux :
$\{a, b, c \} \cup \{a \} = \{a, b, c\}$.

La réunion d'un ensemble et de l'ensemble vide est égale à
cet ensemble : $\{a, b, c \} \cup \emptyset = \{a, b, c \}$.
\end{ex}

\begin{prob}
Démontrer que, si $A \subseteq B$, alors $A \cup B = B$.
\end{prob}

\begin{explain}
On peut aussi considérer une opération «~duale~» de la réunion :
former l'ensemble de tous les !!{element}s qui appartiennent
à la fois à~$A$ et à~$B$. Cette opération s'appelle
l'\emph{intersection} ; elle est représentée dans
la \olref{fig:intersection}.
\begin{figure}
  \olasset{assets/diagrams/intersection.tikz}
  \caption{L'intersection $A \cap B$ de deux ensembles est
    l'ensemble de leurs !!{element}s communs.}
  \ollabel{fig:intersection}
\end{figure}
\end{explain}

\begin{defn}[Intersection]
L'\emph{intersection} de deux ensembles $A$ et $B$, notée
$A \cap B$, est l'ensemble de tous les objets qui sont
des !!{element}s à la fois de $A$ et de~$B$.
\[
A \cap B = \Setabs{x}{x \in A \land x \in B}
\]
Deux ensembles sont dits \emph{disjoints} si leur intersection
est vide, c'est-à-dire s'ils n'ont aucun !!{element} en commun.
\end{defn}

\begin{ex}
Si deux ensembles n'ont aucun !!{element} en commun, leur
intersection est vide. Dans l'exemple suivant, on suppose que
$a$, $b$ et $c$ sont distincts de $0$ et de $1$ :\footnote{Cette
hypothèse, implicite dans l'exemple original, est rendue explicite.}
$\{ a, b, c\} \cap \{ 0, 1\} = \emptyset$.

S'ils ont des !!{element}s en commun, leur intersection est
l'ensemble de ces !!{element}s. En supposant ici $c\neq d$,
on a $\{a, b, c \} \cap \{a, b, d \} = \{a, b\}$.\footnote{Précision
de l'édition française : l'hypothèse $c\neq d$ suffit pour cet exemple.
Sans elle, un élément commun supplémentaire pourrait apparaître.}

L'intersection d'un ensemble et de l'un de ses sous-ensembles
est simplement le plus petit des deux :
$\{a, b, c\} \cap \{a, b\} = \{a, b\}$.

L'intersection de tout ensemble avec l'ensemble vide est vide :
$\{a, b, c \} \cap \emptyset = \emptyset$.
\end{ex}

\begin{prob}
Démontrer rigoureusement que, si $A \subseteq B$, alors
$A \cap B = A$.
\end{prob}

\begin{explain}
On peut aussi former la réunion ou l'intersection de plus de
deux ensembles. Voici une manière commode de traiter le cas général :
réunissons en un seul ensemble tous les ensembles dont nous voulons
prendre la réunion (ou l'intersection). La réunion des ensembles
de départ est alors l'ensemble des objets qui appartiennent à
au moins un !!{element} de cette famille ; leur intersection
est l'ensemble des objets qui appartiennent à chacun de ses
!!{element}s. Pour l'intersection, nous supposons la famille non vide.
\end{explain}

\begin{defn}
Si $A$ est un ensemble d'ensembles, $\bigcup A$ est l'ensemble
des !!{element}s des !!{element}s de~$A$ :
\begin{align*}
\bigcup A & = \Setabs{x}{x \text{ appartient à un !!{element} de } A},
\text{ soit}\\
& = \Setabs{x}{\text{il existe } B \in A
  \text{ tel que } x \in B}
\end{align*}
\end{defn}

\begin{defn}
Si $A$ est un ensemble non vide d'ensembles, $\bigcap A$ est
l'ensemble des objets communs à tous les éléments de~$A$ :
\begin{align*}
\bigcap A & = \Setabs{x}{x \text{ appartient à tout !!{element} de } A},
\text{ soit}\\
 & = \Setabs{x}{\text{pour tout } B \in A, x \in B}
\end{align*}
\end{defn}
\noindent\emph{Précision de l'édition française.}
L'hypothèse $A\neq\emptyset$ est nécessaire ici : sans elle, la
condition universelle serait vérifiée par tout objet. Si l'on travaille
dans un ensemble ambiant fixé $U$, on peut en revanche définir
l'intersection de la famille vide comme étant $U$.

\begin{ex}
Supposons que $A = \{ \{ a, b \}, \{ a, d, e \}, \{ a, d \} \}$,
avec $b\neq d$.\footnote{Précision de l'édition française :
l'hypothèse $b\neq d$, implicite dans l'exemple original,
garantit que les trois ensembles n'ont aucun élément commun autre que $a$.}
Alors $\bigcup A = \{ a, b, d, e \}$ et $\bigcap A = \{ a \}$.
\end{ex}
\begin{prob}
	Montrer que, si $A$ est un ensemble et $A \in B$,
	alors $A \subseteq \bigcup B$.
\end{prob}

On peut procéder de même pour une suite d'ensembles
$A_1$, $A_2$, \dots
\begin{align*}
\bigcup_i A_i & = \Setabs{x}{x \text{ appartient à l'un des } A_i}\\
\bigcap_i A_i & = \Setabs{x}{x \text{ appartient à chacun des } A_i}.
\end{align*}

Si les ensembles sont \emph{indexés} par un ensemble $I$,
c'est-à-dire si l'on considère un ensemble $A_i$ pour chaque
$i \in I$, on peut aussi employer les abréviations suivantes,
avec $I\neq\emptyset$ pour l'intersection sans ensemble ambiant fixé :
\begin{align*}
	\bigcup_{i \in I} A_i & = \bigcup \Setabs{A_i }{i \in I}\\
	\bigcap_{i \in I} A_i & = \bigcap\Setabs{A_i}{i \in I}
\end{align*}

Enfin, on peut considérer l'ensemble de tous les !!{element}s
de~$A$ qui n'appartiennent pas à~$B$. La \olref{difference}
en donne une représentation.

\begin{figure}
  \olasset{assets/diagrams/difference.tikz}
  \caption{La différence $A \setminus B$ de deux ensembles
    est l'ensemble des !!{element}s de~$A$ qui ne sont pas
    des !!{element}s de~$B$.}
  \ollabel{difference}
\end{figure}

\begin{defn}[Différence]
La \emph{différence ensembliste}~$A \setminus B$ est l'ensemble
de tous les !!{element}s de $A$ qui ne sont pas des !!{element}s
de~$B$, c'est-à-dire
\[
A\setminus B = \Setabs{x}{x\in A \text{ et } x \notin B}.
\]
\end{defn}

\begin{prob}
	Démontrer que, si $A \subsetneq B$,
	alors $B \setminus A \neq \emptyset$.
\end{prob}

\end{document}
